\documentclass[../../../uem_mat_tok_q.tex]{subfiles}

\begin{document}
    Oを原点とする座標空間において，不等式 $|x|\leqq 1$，$|y|\leqq 1$，$|z|\leqq 1$ の%
    表す立方体を考える．その立方体の表面のうち，$z<1$ を満たす部分を$S$とする．

    以下，座標空間内の2点A，Bが一致するとき，線分ABは点Aを表すものとし，%
    その長さを0と定める．

    \begin{enumerate}
        \item 座標空間内の点Pが次の条件\,\ref{uem_mat_tok_2023_f_sci_06_q:1}，%
            \ref{uem_mat_tok_2023_f_sci_06_q:2}\,をともに満たすとき，%
            点Pが動きうる範囲$V$の体積を求めよ．
            \begin{enumerate}
                \item $\mathrm{OP}\leqq\sqrt{3}$
                    \label{uem_mat_tok_2023_f_sci_06_q:1}
                \item 線分OPと$S$は，共有点を持たないか，点Pのみを共有点に持つ．
                    \label{uem_mat_tok_2023_f_sci_06_q:2}
            \end{enumerate}
        \item 座標空間内の点Nと点Pが次の条件\,\ref{uem_mat_tok_2023_f_sci_06_q:3}，%
            \ref{uem_mat_tok_2023_f_sci_06_q:4}，\ref{uem_mat_tok_2023_f_sci_06_q:5}\,をすべて満たすとき，%
            点Pが動きうる範囲$W$の体積を求めよ．必要ならば，%
            $\sin\alpha=\myfraca{1}{\sqrt{3}}$ を満たす%
            実数$\alpha\,\left(0<\alpha<\myfraca{\pi}{2}\right)$を用いてよい．
            \begin{enumerate}[resume]
                \item $\mathrm{ON}+\mathrm{NP}\leqq\sqrt{3}$
                    \label{uem_mat_tok_2023_f_sci_06_q:3}
                \item 線分ONと$S$は共有点を持たない．
                    \label{uem_mat_tok_2023_f_sci_06_q:4}
                \item 線分NPと$S$は，共有点を持たないか，点Pのみを共有点に持つ．
                    \label{uem_mat_tok_2023_f_sci_06_q:5}
            \end{enumerate}
    \end{enumerate}
\end{document}